\documentclass[11pt,a4paper]{article}

\usepackage[utf8]{inputenc}
\usepackage[margin=1in]{geometry}
\usepackage{amsmath,amssymb,amsfonts}
\usepackage{booktabs}
\usepackage{tabularx}
\usepackage{graphicx}
\usepackage{natbib}
\usepackage{hyperref}
\usepackage{microtype}
\usepackage{setspace}

\hypersetup{
    colorlinks=true,
    linkcolor=blue!70!black,
    citecolor=blue!70!black,
    urlcolor=blue!70!black
}

\title{\textbf{Central Money, Decentralized Credit Prices: Allocation, Shock Transmission, and Liquidity Backstops in a Calibrated Agent-Based Economy}}

\author{
    \textbf{Manuscript Draft} \\
    \textit{Prepared for Journal of Economic Dynamics and Control}
}

\date{\today}

\begin{document}

\maketitle

\begin{abstract}
Can retail credit prices be decentralized while the currency and settlement asset remain centralized? We study this institutional separation in a stock-flow-consistent, agent-based economy with heterogeneous firms, competing banks, a common unit of account, interbank settlement, emergency liquidity, and explicit bank resolution. The treatment changes the sensitivity of loan quotes to borrower and bank state rather than imposing a mechanical average-rate gap. Bank populations and financial parameters are disciplined using US bank and small-business lending data from 2022--2024, with 2025 used for out-of-sample validation subject to a disclosed model-selection limitation. In 809 matched confirmatory seeds, market pricing raises credit-weighted borrower productivity by 0.075 model units but raises the unfunded-demand share by 0.493. Thus, stronger local risk pricing improves selection on the intensive margin while substantially contracting the extensive margin. Following a positive demand shock, the market regime reduces cumulative new-credit and output responses by 0.145 and 0.351 model units, respectively. Reserve abundance and facility design materially alter relative credit and unresolved-liquidity outcomes. Ablations and local sensitivity identify borrower-risk pricing as the central mechanism; global stress tests show that the shock-transmission result is conditional, not universal. A secondary DeepSeek R1 8B exercise completes 30 matched pairs with 900 valid calls and broadly reproduces higher market-state loan rates. The results characterize an institutional tradeoff, not an optimal monetary-policy ranking.
\vspace{0.8em}

\noindent \textbf{Keywords:} Agent-based computational economics; Credit rationing; Market-determined interest rates; Liquidity backstops; Stock-flow consistency. \\
\textbf{JEL Classification:} E43, E51, E58, G21, C63.
\end{abstract}

\onehalfspacing

\section{Introduction}

Debates over market-determined interest rates often bundle together several institutions that need not move together. The unit of account may be public; the settlement asset may be issued by a central monetary authority; commercial banks may create deposit liabilities through lending; and the price and allocation of an individual loan may nevertheless respond to local borrower and bank information. Treating this entire arrangement as either ``centralized'' or ``decentralized'' hides the mechanism of interest. This paper separates the centralization of money and settlement from the decentralization of retail credit pricing.

We ask three narrower questions:
\begin{enumerate}
    \item When loan quotes respond more strongly to borrower leverage and bank funding conditions, does credit move toward more productive firms, and how much additional demand goes unfunded?
    \item How does that pricing institution change the propagation of a positive aggregate demand shock?
    \item How do reserve abundance and emergency-liquidity design condition the answer?
\end{enumerate}
These questions concern pricing, allocation, settlement, and institutional dependence. They do not ask whether a model-generated rate is the optimal policy rate for the United States.

The model is a controlled computational laboratory. Thirty heterogeneous firms interact with five banks over 24 quarterly periods in the confirmatory design. Firms request loans with a quantity, maturity, purpose, expected return, and maximum acceptable rate. Banks return partial offers containing an approved quantity and a complete pricing decomposition. Firms accept compatible offers from cheapest upward subject to a lender-concentration limit. Credit funds both working capital and productive investment. Production, sales, inventories, loan servicing, payment settlement, liquidity support, default, recovery, and bank resolution change explicit balance sheets. The monetary authority issues the settlement asset and can provide emergency liquidity or resolution capital, but it does not select the borrower-specific market-regime loan quote.

The comparison has two regimes. Under administered pricing, the benchmark is a policy rate plus a calibrated spread and only one quarter of local funding, risk, liquidity, capital, and inflation components passes into the quote. Under market pricing, the benchmark is the required real return plus a calibrated intercept and local components pass through fully. The intercepts are aligned so that the experiment does not mechanically encode a mean-rate treatment. The economic treatment is the information sensitivity of the price.

The paper makes four contributions:
\begin{enumerate}
    \item \textbf{Institutional decomposition:} Centralized money is compatible with decentralized loan-price formation.
    \item \textbf{Joint allocation and quantity measurement:} It measures the joint allocation and quantity consequences of that decomposition using requested, offered, accepted, and unfunded principal rather than treating a rate coefficient as the economic result.
    \item \textbf{Interaction with settlement architecture:} It embeds the credit market in explicit production and settlement blocks and estimates how reserves and backstops interact with the pricing institution.
    \item \textbf{Strict evidence workflow:} It applies an unusually strict evidence workflow for a computational model: empirical source hashes, deterministic calibration, a disclosed validation gate, pilot-based power, independent matched seeds, immutable event ledgers, exact consolidation checks, and generated paper assets.
\end{enumerate}

The main result is a tradeoff. Market pricing raises credit-weighted borrower productivity by 0.0751, but raises the average unfunded-demand share by 0.4932. The market regime therefore selects more productive recipients among firms that receive credit while denying much more requested credit overall. It is not correct to call this an unqualified efficiency improvement. Following the positive demand shock, market pricing produces a 0.1445 smaller cumulative credit impulse and a 0.3512 smaller output impulse. These effects are precisely estimated in the confirmatory design but small relative to the total positive response. Borrower-risk pricing is the central mechanism: removing it nearly eliminates the selection and shock-dampening contrasts.

Institutional support matters. At low reserves with no emergency facility, the two regimes differ in both cumulative credit and unresolved settlement shortfalls. High reserves or a penalty facility remove unresolved shortfalls at the declared anchor cells, while the relative credit effect persists. The proper conclusion is not that scarce liquidity is desirable. It is that a credit-pricing institution cannot be evaluated independently of the settlement and backstop architecture supporting it.

Robustness results discipline the scope of these conclusions. Population and bank-count variations preserve the allocation and rationing results. A corrected, separately frozen concentration addendum preserves the selection and rationing effects under low and high opening deposit concentration. A $\pm 10\%$ local sensitivity design preserves the demand-shock dampening signs in 12 of 13 parameter sets; a modest reduction in the borrower-risk coefficient moves both estimates slightly above zero. An extreme global stress map contains many sign reversals. The calibrated shock-transmission finding is therefore a local institutional result, not a theorem about all possible economies.

The rest of the paper presents the literature and institutional contribution, defines the model, describes data and calibration, states the experimental design, reports empirical results across allocation, shock transmission, and liquidity interactions, and discusses robustness, the bounded LLM exercise, limitations, and implications.

\section{Related Literature and Contribution}

The paper connects credit rationing, bank-mediated macroeconomic models, agent-based computational economics, monetary implementation, and emerging LLM agent methods. \citet{stiglitz1981credit} established that loan-market clearing need not occur solely through an interest rate when borrower risk and information are imperfect. Our object differs from the canonical theoretical problem: borrower state is observable to the programmed bank, and the experiment varies how strongly that state enters the quote. Rationing remains central because a higher quote can exceed expected project return and because bank capital and lender concentration limit quantities.

Multi-bank artificial credit-market models study relationship lending, supervision, and quantity allocation \citep{zhang2018credit, neuner2025agent}. Macro-financial agent-based models link bank balance sheets to firm production and aggregate fluctuations \citep{ashraf2017banks, caiani2016agent, popoyan2017taming}. Our closest point of departure is this decentralized bank-firm matching tradition. The increment is to isolate the rate institution while retaining a centralized currency and settlement system, log the complete quote decomposition, and predeclare allocation and rationing as primary outcomes.

The study also belongs to the computational-laboratory view of agent-based economics \citep{dosi2017macroeconomic, hendry2025agent}. Heterogeneous state and local interaction produce aggregate outcomes that generally lack a convenient closed form. This flexibility creates a corresponding burden: the researcher must show that results are not accidents of code, seed choice, hidden accounting residuals, or post-result tuning. We therefore treat independent simulation seeds as the sampling unit, separate pilot and confirmatory namespaces, record every event in normalized tables, and publish adverse and sign-reversing robustness results \citep{detail2026bankofengland}.

Finally, recent work uses language models as economic agents. We do not make LLM behavior the identification strategy. A small local DeepSeek model is used only to ask whether bounded decisions under identical institutional states display qualitatively comparable regime differences. The core claims continue to rest on transparent rule agents, for which every behavioral coefficient and mechanism switch is inspectable.

\section{The Model}

\subsection{Agents, Timing, and Accounting}

The main economy contains firms indexed by $i \in \{1,\dots,30\}$, banks indexed by $b \in \{1,\dots,5\}$, one aggregate household/capital-goods supplier, and a monetary authority. Wages and the goods price are normalized to one. The horizon is 24 periods. A period proceeds through pending bank resolution, production, incumbent-loan servicing, new-credit clearing, goods-market clearing, experimental-loan servicing, deposit-funding reallocation, settlement-liquidity clearing, reserve-buffer management, reserve remuneration, insolvency marking, and event recording.

The accounting identities are checked after every period. Household plus firm deposits equal bank deposit liabilities. Firm debt equals bank customer loans. Interbank assets and liabilities, authority money creation, emergency lending, and resolution injections are separately recorded. These are executable model conditions: a violation fails the run rather than becoming an unexplained residual.

\subsection{Production, Demand, and Firm Finance}

Firm $i$'s output at period $t$ is given by a Cobb-Douglas technology:
\begin{equation}
Y_{it} = A_i S_t K_{it}^{\alpha} L_{it}^{1-\alpha},
\end{equation}
where $A_i$ is firm-specific productivity, $S_t$ is a common productivity-shock multiplier, $K_{it}$ is productive capital, and $L_{it}$ is labor financed from deposits and available working capital. Capital depreciates before production at rate $\delta$. New accepted credit is divided between working capital and investment using a calibrated or declared investment share. Investment is a deposit transfer to the household/capital-goods supplier and an equal addition to firm real capital; it is not free net worth.

Expected project return equals the common expected-return parameter multiplied by firm productivity and the current demand-shock multiplier. Desired credit is base demand plus a return-gap component, again multiplied by the demand shock, and capped by a leverage constraint. Applications carry the desired principal, 19-period calibrated maturity, mixed purpose, expected return, and maximum acceptable rate.

Household planned consumption is bounded by deposits and depends on current labor/capital income and lagged deposit wealth. Actual consumption is the minimum of planned expenditure, available goods, and household deposits. Sales are allocated in proportion to each firm's available inventory. Unsold output and unmet consumption remain explicit.

\subsection{Loan Pricing and Clearing}

Let $z_{ibt}$ collect local funding, borrower-risk, liquidity, capital, and inflation components for firm $i$ and bank $b$. The offered nominal rate is:
\begin{equation}
r_{ibt} = q_j + s_j + \lambda_j z_{ibt} + \eta_{ib},
\end{equation}
where $j \in \{A, M\}$ denotes administered ($A$) or market ($M$) pricing, and $\eta_{ib}$ is a fixed relationship component generated from the matching seed. More explicitly:
\begin{equation}
z_{ibt} = f_{bt} + \gamma_{\text{risk}}\max(0, \ell_{it}) + \gamma_{\text{liq}}\max(0, \bar{m} - m_{bt}) + \gamma_{\text{cap}}\max(0, \bar{k} - k_{bt}) + \gamma_{\pi}\max(0, \pi^e_{bt}).
\end{equation}
The administered benchmark $q_A$ is the policy rate, $s_A$ is the calibrated administered spread, and $\lambda_A = 0.25$. The market benchmark $q_M$ is the required real return, $s_M$ is the calibrated market intercept, and $\lambda_M = 1.0$. The funding term increases when loans exceed the deposit target or the bank uses costly reserve/interbank funding. Intercepts are calibrated to avoid an imposed baseline mean-rate gap.

An active bank offers up to the lesser of requested credit and its residual regulatory-capital capacity. It rejects when inactive, already in unresolved liquidity failure, out of regulatory capacity, or when its quote exceeds the firm's expected return. The firm sorts feasible offers by price and accepts from cheapest upward until demand is filled or the lender-concentration limit binds. This produces direct measures of requested, approved, accepted, and unfunded quantities.

\subsection{Loans, Liquidity, and Resolution}

Experimental loans amortize linearly. Interest uses the annual nominal rate converted to the quarterly model frequency. If deposits cannot meet scheduled principal and interest, the loan defaults. Recovery first records cash/deposit payment and then a fixed fraction of remaining productive capital; the residual is written off against bank profit. Incumbent empirical loan stocks are assigned as firm liabilities and rolled using calibrated growth rates, avoiding an unbacked opening bank asset.

Payments are final. Cross-bank transfers change reserves. At the end of the period, banks below the settlement floor borrow first in the interbank market. Remaining need may be met by an unavailable, limited, or penalty emergency facility depending on the scenario. Any residual is an unresolved liquidity shortfall; earlier payments are not reversed.

A bank with negative equity enters resolution and cannot make new offers. After one period, the authority injects the amount needed to restore the regulatory capital ratio. The injection increases reserves, equity, base money, and a reported resolution-cost variable. Emergency liquidity and solvency resolution are therefore distinct contracts.

\section{Data and Calibration}

Bank balance-sheet and income information comes from public FDIC Call Report bulk files. Business-loan rates, amounts, maturities, standards, applicant quality, and demand use the Federal Reserve FR 2028D Small Business Lending Survey, with Financial Accounts and H.8/FRED series used as aggregate checks. Nominal loan amounts are converted to constant 2025Q4 prices for empirical processing; the simulation itself uses normalized model units.

The estimation sample covers 2022Q1--2024Q4. Eligible observations are active, domestically chartered, FDIC-insured commercial banks with positive assets, deposits, loans, and equity. Impossible values are removed before 1st/99th percentile winsorization within quarter and size stratum. The pipeline preserves equal-bank and asset-weighted moments, logs exclusions and missingness, and fails on checksum changes, missing quarters, duplicate bank-quarter keys, accounting impossibilities, or unrecognized field drift.

Bank populations are initialized by joint empirical resampling of assets, deposits, loans, liquid assets, equity, C\&I share, growth, and charge-off state. The median sampled bank's deposits are normalized to 100 model units while ratios and cross-sectional dependence are preserved. We target capital/assets, loans/deposits, liquidity/deposits, C\&I share, deposit and loan growth, charge-offs, loan amount and maturity, and effective-rate/spread moments.

Calibration minimizes variance-standardized moment errors using 200 bank-cluster bootstrap draws and ten recorded differential-evolution starts. The accepted bundle has normalized RMSE 0.161, below the 0.25 gate. Seven of eight 2025 validation groups fall within empirical 95\% bootstrap intervals and none is more than two standard errors away.

This validation evidence has an important qualification. A preserved candidate first placed six of eight groups inside their intervals. The accepted bundle was written after that result was inspected and includes a slightly revised administered spread plus additional transition parameters. We therefore call 2025 an out-of-sample validation period used with an acceptance gate, not a pristine untouched holdout. A later frozen data vintage would be required for a fully independent external test.

\section{Experimental Design and Empirical Strategy}

The empirical evaluation focuses on three primary institutional mechanisms:
\begin{enumerate}
    \item \textbf{Quote Pass-Through and Credit Allocation:} A manipulation check confirms that market pricing delivers stronger local-state pass-through, followed by an evaluation of changes in credit-weighted productivity and the unfunded-demand share.
    \item \textbf{Demand-Shock Transmission:} Under the recent-US calibration, we test whether market pricing dampens the cumulative new-credit and output response to a positive aggregate demand shock.
    \item \textbf{Liquidity and Emergency Backstops:} We examine how reserve abundance and facility design condition relative credit volume and settlement liquidity shortfalls across regimes.
\end{enumerate}

The confirmatory main design includes baseline and positive demand-shock scenarios, two regimes, and 809 matched seeds. The shock begins in period 8, lasts four periods, and raises demand by 25 percent. Cumulative responses use periods 8--23. The liquidity-backstop experiment crosses five reserve anchors, four facility settings, two regimes, and 40 matched seeds. Mechanism ablations, topology cells, and 100 Latin-hypercube stress sets are robustness exercises. The frozen design totals 8,096 rule runs.

Inference is across independent seeds, never agent-period rows. Allocation estimands use matched market-minus-administered seed differences. Demand-shock transmission effects use a matched difference-in-differences (market shock minus market baseline, less administered shock minus administered baseline). Liquidity and facility interactions use predeclared interaction contrasts. Confidence intervals use 10,000 percentile bootstrap draws resampling matched seeds. Holm correction is applied within each two-outcome primary family.

\section{Results}

\subsection{Credit Allocation and the Selection--Rationing Tradeoff}

The manipulation check succeeds. Full rather than attenuated pass-through raises the local-state pass-through coefficient by 0.75. The slope of quoted rates on borrower leverage is 0.02145 higher in the market regime (95\% interval [0.02128, 0.02163]). This confirms the treatment but is not itself a welfare result.

Market pricing raises credit-weighted borrower productivity by 0.07511 among 804 matched pairs with positive credit in both regimes (95\% interval [0.07206, 0.07821], Holm-adjusted $p \approx 0.0002$). Mean levels are 1.07546 under administered pricing and 1.15104 under market pricing. Five of 809 market runs originate no new credit, making a credit-weighted productivity statistic undefined; those observations are disclosed rather than assigned artificial values.

The allocation gain comes with severe quantity contraction. The mean unfunded-demand share rises by 0.49318 (95\% interval [0.48349, 0.50281], Holm-adjusted $p \approx 0.0002$), from 0.03812 under administered pricing to 0.53130 under market pricing. Quoted-rate dispersion rises by 0.00496 and the leverage gradient of rationing rises by 0.39114. This documents a selection-versus-quantity tradeoff: credit recipients are more productive, but far less requested credit is funded.

Baseline defaults and write-offs are zero in both regimes. The experiment cannot identify relative default performance in the calibrated baseline. Output per unit of new credit is mechanically high in the market regime because new credit can approach zero, so it remains secondary and is not interpreted as a welfare statistic.

\subsection{Macroeconomic Demand-Shock Transmission}

The positive demand shock raises cumulative new credit by 4.00691 under administered pricing and 3.86237 under market pricing. The matched difference-in-differences is $-0.14454$ (95\% interval [$-0.21664, -0.07107$], Holm-adjusted $p \approx 0.0003$). This is approximately 3.6 percent of the administered credit impulse.

The administered output impulse is 30.86338, compared with 30.51219 under market pricing. The difference is $-0.35119$ (95\% interval [$-0.45484, -0.24840$], Holm-adjusted $p \approx 0.0002$), approximately 1.1 percent of the administered output impulse. The hypothesis is supported at the calibration: stronger local pricing modestly dampens the positive shock rather than reversing it.

The borrower-risk-pricing ablation nearly eliminates both shock transmission effects. In the corrected local sensitivity design, 12 of 13 $\pm 10\%$ parameter sets retain negative credit and output signs. Reducing risk pricing by 10 percent produces small positive estimates with intervals spanning zero; raising it strengthens dampening. The global stress map contains many reversals. These results identify the mechanism and the boundary simultaneously: demand-shock dampening is locally stable across most examined parameters but depends on sufficiently strong borrower-risk pricing.

\subsection{Settlement Liquidity and Emergency Facility Interactions}

At low reserves with an unavailable facility, market pricing produces 8.59854 less cumulative credit than administered pricing and 0.73728 less unresolved liquidity shortfall. At high reserves without a facility, unresolved shortfalls are zero in both regimes and the market credit gap is $-12.25852$. At low reserves with a penalty facility, shortfalls are again zero and the credit gap is $-11.73893$.

Moving from high to low reserves when the facility is unavailable changes the market-versus-administered credit contrast by $+3.65998$ and its liquidity contrast by $-0.73728$. Replacing a penalty facility with no facility at low reserves changes the credit contrast by $+3.14039$ and the liquidity contrast by $-0.73728$. All four declared intervals exclude zero after within-comparison Holm correction.

The liquidity sign must not be read as a general advantage of the market regime. It occurs in the anchor where unresolved shortfalls exist, and it partly reflects the market regime's smaller credit volume. The key liquidity finding is that reserve and facility settings materially condition the relative regime effect.

\subsection{Topology, Concentration, and LLM Robustness}

The allocation and rationing signs survive 30-firm/three-bank, 30/5, and 100/5 population checks. The original concentration cells were invalid because empirical initialization bypassed the switch. They remain archived. A separately frozen 150-run correction raises mean opening deposit HHI from 0.2012 in the low-concentration 30/5 cell to 0.2906 in the high cell. Credit-weighted productivity effects are 0.0659 and 0.0647, while unfunded-share effects are 0.6166 and 0.5970, respectively. These are within-cell regime checks, not causal estimates of concentration.

DeepSeek R1 8B completes 30 matched pairs, with all 900 main calls valid and no retries. The market prompt raises the average offered rate by 0.00552 (95\% interval [0.00383, 0.00732]) and lowers approved principal by 0.42967 (95\% interval [$-0.69898, -0.15097$]). Requested principal differs by $-0.43347$ with an interval spanning zero. The exercise supports the qualitative rate response under bounded actions, but one small model and three prompt templates do not establish behavioral external validity.

\section{Discussion and Limitations}

The findings reject a simple winner-versus-loser framing. Full local-state pricing directs a smaller pool of credit toward more productive firms. Whether that tradeoff is socially preferable depends on welfare weights, missing household heterogeneity, long-run entry and exit, distributional incidence, and the value of projects left unfunded. The current design does not supply those objects and therefore does not claim optimality.

Several limitations are substantive:
\begin{itemize}
    \item The household sector is aggregate; wages and the price level are normalized to one.
    \item The horizon is 24 quarters, which abstracts from multi-decade secular cycles.
    \item Production parameters are calibrated rather than directly identified by bank-level microdata.
    \item Baseline defaults are too rare to compare realized credit losses empirically across steady states.
    \item The resolution authority recapitalizes banks transparently but abstracts from legal resolution regimes and panics.
    \item The 2025 validation period contains model-selection leakage as disclosed.
    \item Global sensitivity produces sign reversals, and local sensitivity identifies borrower-risk pricing as a nearby boundary.
    \item The LLM exercise is small and prompt-conditioned.
\end{itemize}

These limits sharpen rather than erase the contribution. The evidence supports a precise conditional statement: in this calibrated institutional model, stronger borrower- and bank-state pass-through changes who receives credit, substantially increases rationing, slightly dampens a positive demand shock, and interacts with settlement liquidity and backstop design.

\section{Conclusion}

Centralized money does not logically require centrally administered retail credit prices. In the simulated economy studied here, decentralizing the local information content of loan quotes produces a measurable selection-versus-quantity tradeoff. Funded credit shifts toward more productive borrowers while roughly half more requested credit remains unfunded. Demand-shock transmission is modestly dampened at the recent-US calibration, and the result depends on borrower-risk pricing. Reserve abundance and emergency facilities materially change relative credit and liquidity outcomes.

The policy lesson is institutional rather than prescriptive. Credit-pricing rules, settlement assets, prudential capacity, and liquidity backstops form a joint system. Changing one element while holding the others implicit can misstate the mechanism. Future work should add heterogeneous households, endogenous entry and bank competition, a longer default cycle, empirically identified production parameters, and a truly untouched validation vintage.

\bibliographystyle{elsarticle-harv}
\bibliography{references}

\end{document}
